% hpsatviews --- Earth Science Informatics (Springer)
% Compilar: pdflatex -> bibtex -> pdflatex x2
%
% TIPO DE ARTICULO: methodology article, subtipo computational. Hay que
% declararlo en la interfaz de envio, y es lo que fija la estructura de la
% seccion de metodos en las subsecciones Algorithm / Implementation / Testing
% que sigue este archivo.
%
% CLASE: se compila con `article` para poder trabajar en cualquier maquina.
% Para el envio final Springer recomienda su propia plantilla, que se instala
% con
%     sudo apt install texlive-publishers
% y entonces basta cambiar la linea de \documentclass por
%     \documentclass[sn-mathphys-num]{sn-jnl}
% mas el bloque de portada. El cuerpo no usa ninguna macro propia de la clase,
% asi que el cambio no lo toca.
%
% ESTE NO ES EL MANUSCRITO DE SoftwareX REFORMATEADO, y tampoco el de JTECH.
% ESIN exige que la contribucion principal sea a la informatica, y su alcance
% excluye por nombre los articulos centrados en percepcion remota o
% meteorologia como dominio de aplicacion. Prueba para cada parrafo: esto le
% sirve a alguien que NO procese datos GOES? Si no, sobra.
%
% Plan completo y pendientes: paper/plan_esin.md

\documentclass[11pt]{article}

\usepackage[utf8]{inputenc}
\usepackage[T1]{fontenc}
\usepackage{booktabs}
\usepackage{natbib}

% TITULO. El actual viene de JTECH y lidera con lo operativo y con GOES, que
% es justo lo que la advertencia de alcance de ESIN castiga. Alternativas que
% ponen la informatica al frente, a decidir junto con el abstract:
%  - Quadratic chunk-index cost in chunked scientific storage: characterization,
%    removal, and a reproducible benchmarking protocol
%  - Where the time goes in geostationary product pipelines: chunk-index cost,
%    device-resident processing, and a measured protocol
%  - Characterizing I/O and precision bottlenecks in chunked-HDF5 Earth
%    observation pipelines
\title{Where the time goes in operational GOES-R ABI product generation:
parallel decoding, device-resident processing, and a measured comparison}

% TODO decidir coautores. Pendiente desde el resumen de la RAUGM. Bloquea
% tambien el bloque Author contributions del cierre. Springer pide ORCID de
% 16 digitos de cada autor y correo del autor de correspondencia.
\author{Alejandro Aguilar Sierra\\
\small Laboratorio Nacional de Observaci\'on de la Tierra (LANOT),\\
\small Universidad Nacional Aut\'onoma de M\'exico, Mexico City, Mexico\\
\small \texttt{asierra@unam.mx}}

\date{}

\begin{document}

\maketitle

% Entre 150 y 250 palabras; el actual tiene 245, asi que no hay margen para
% crecer. Reescribir AL FINAL sobre la tesis informatica, no sobre el encuadre
% operativo heredado de JTECH. La revista prohibe abreviaturas sin definir: hoy
% usa ABI y GPU sin expandir.
\begin{abstract}
Operational monitoring with geostationary imagery is bounded by a
deadline: the Advanced Baseline Imager delivers a new full disk every ten to
fifteen minutes, and the interval between file arrival and image availability
determines whether a product informs a decision. The algorithms defining
standard GOES-R composites are settled, but the per-scene cost of applying them
at operational resolution is uncharacterized. We ask where that time is
spent and how much is recoverable without altering them. Decoding
dominates, and much of its cost is hidden in retrieving the HDF5 chunk index,
which grows quadratically with the chunk count when queried per chunk and is
undetectable at coarser resolutions. Retrieving that index in a single
traversal, inflating the raw chunks concurrently, omitting an overview pyramid
downstream stages discard, and executing the per-pixel arithmetic as one
device-resident chain reduce the wall time of a full-disk true-color composite
from $30.6$~s to $9.25$~s on the same processors as the reference tool, and to
$3.02$~s on a GPU, with outputs agreeing to a mean absolute difference of
$1.5\%$ of full scale. Whether GPU offloading pays proves to be a property of
the device, not of the approach: the geolocation stages require double
precision and run $15.4\times$ slower on one accelerator than another, while
the single-precision stages differ by $1.9\times$, and at full resolution the
working set exceeds the smaller card entirely. We also report three
measurement failure modes, two of which would have overstated our results, and
the protocol avoiding them.
\end{abstract}

% TODO elegir de 4 a 6, apuntando a informatica y no a percepcion remota.
% Candidatas: chunked storage, HDF5, scientific data I/O, GPU computing,
% reproducible benchmarking, geostationary satellite data.
\textbf{Keywords:} \textit{(pendiente)}

% El "significance statement" es un requisito de AMS que ESIN no pide. Se
% conserva porque sirve casi tal cual para la carta de presentacion:
%
%% When a hurricane intensifies or a volcano erupts, the satellite image behind
%% the warning is useful only if it arrives quickly. Producing one is routine;
%% producing it fast is not, and the time it takes has rarely been examined. We
%% measured where that time goes for a standard product and found most of it not
%% in the calculations but in reading the compressed data, where one step slows
%% disproportionately as resolution rises. Correcting that and rearranging the
%% remaining work cuts the delay threefold without changing the imagery at all.
%% Graphics accelerators help or not depending on the card, and we describe how to
%% measure this reliably, since it is easy to get wrong.

\section{Introduction}
\label{sec:intro}

The Advanced Baseline Imager aboard the GOES-R series delivers a new full disk
every ten to fifteen minutes, and a new contiguous-United States sector every
five. Each full-disk time stamp comprises sixteen spectral bands at resolutions
down to $0.5$~km---roughly $0.94$~GB of compressed NetCDF for a representative
scene---which must be decoded, radiometrically corrected, normalized, composited
and often reprojected before an interpretable image exists. For a forecast
office or a hazard-monitoring center, the interval between the arrival of those
files and the availability of the image is an operational quantity: a volcanic
ash product or a rapid-intensification sequence that takes several minutes to
render has consumed part of the window it was meant to inform. At our
laboratory these products are distributed to the national civil-protection and
naval agencies responsible for early warning, and the requirement is not peak
throughput but short and, above all, predictable turnaround.

The algorithms that define these products are settled and should not be
reinvented. The synthesis of a green band absent from the instrument
\citep{Bah2018,Miller2016}, the Rayleigh correction and its lookup-table
formulation \citep{Bucholtz1995,Bodhaine1999,Scheirer2018,PySpectralLUT2018},
and the contrast transformations that give the familiar composites their
appearance are documented and implemented in mature community software, of
which \texttt{satpy} and \texttt{geo2grid} \citep{Geo2Grid2026} are the
reference. Those packages are built for the breadth of research use---many
instruments, several resampling backends, an exploratory workflow over a
general-purpose array stack---and that breadth is precisely what an operational
pipeline does not need and does pay for.

What is missing is not another implementation but a measurement. Performance
claims about satellite-processing software are commonly stated without
specifying what was compared, on which hardware, or against which
configuration of the reference tool, and are consequently neither verifiable
nor portable. This paper asks a narrower and answerable question: for a
standard GOES-R composite at operational resolution, where does the per-scene
time actually go, and how much of it is recoverable without altering the
algorithms? We answer it by building an engine designed for that regime and
then measuring it against the established tool under a protocol we state
explicitly.

Our contributions are the following.

\begin{enumerate}
\item We identify a quadratic cost in the retrieval of the HDF5 chunk index
that dominates decoding at full-disk resolution and is invisible at coarser
scales, and show that a single-traversal alternative removes it
(Section~\ref{sec:decoding}). The finding applies to any consumer of chunked
NetCDF-4, independently of the software presented here.

\item We describe a device-resident organization of the processing chain, in
which the transfer cost is paid once per scene rather than once per stage, and
quantify the memory it demands (Sections~\ref{sec:device}
and~\ref{sec:portability}).

\item We report three failure modes of performance measurement that we
encountered, each of which yields a plausible and reproducible number that
means something other than what it appears to, and state the protocol that
avoids them (Section~\ref{sec:methodology}). Two of the three would have
overstated our own results.

\item We measure the resulting engine against \texttt{geo2grid} on identical
hardware, scene and product, and verify that the two produce equivalent output
before comparing their times (Section~\ref{sec:results}).

\item We show, with per-stage measurements on two accelerators, that whether
GPU offloading pays for this workload is governed by the device's
double-precision throughput and memory capacity rather than by the approach
itself, and that on one widely deployed card the accelerated path is
unavailable at full resolution (Section~\ref{sec:portability}).
\end{enumerate}

Following \citet{LiraChavez2010} we distinguish the \emph{image}, the numerical
representation of the observed scene, from the \emph{view} derived from it for
human interpretation, and from the \emph{product} that a view becomes once it
corresponds to a concept recognized by the community. The work reported here
concerns the transformation from the first to the third, and makes no claim on
the retrieval algorithms that precede it.

%%%%%%%%%%%%%%%%%%%%%%%%%%%%%%%%%%%%%%%%%%%%%%%%%%%%%%%%%%%%%%%%%%%%%
\section{Methods}
\label{sec:methods}

\subsection{Algorithm: products and reference definitions}
\label{sec:algorithms}

%% TODO. Breve y con propósito: fijar la línea base de equivalencia, para que
%% la sección de Resultados pueda afirmar que se compara el MISMO producto.
%%  - Verde sintético: G = 0.465 B + 0.465 R + 0.07 NIR (coeficientes CIMSS).
%%  - Rayleigh por LUT derivada de pyspectral; relajación por nubes.
%%  - Estiramiento por tramos y realce por razón, los de geo2grid.
%% Insistir: los algoritmos son los publicados. Lo que cambia es el sustrato.

%%%%%%%%%%%%%%%%%%%%%%%%%%%%%%%%%%%%%%%%%%%%%%%%%%%%%%%%%%%%%%%%%%%%%
\subsection{Implementation}
\label{sec:implementation}

\subsubsection{Parallel decoding of the input}
\label{sec:decoding}

A GOES-R ABI full disk arrives as sixteen separately compressed NetCDF-4 files.
Decoding them is the first stage of any product chain and, once the arithmetic
is parallelized, the one that dominates. The obstacle is not the volume of data
but the interface: reading a variable through the NetCDF or HDF5 API routes
every chunk through the library's filter pipeline, which is single-threaded. On
a $0.5$~km band---$21{,}696 \times 21{,}696$ samples, or $4.7 \times 10^{8}$
values---that pipeline leaves the remaining cores idle while one thread
inflates the whole variable.

The alternative is to bypass the filter pipeline: retrieve each chunk in its
stored, still-compressed form, then inflate the chunks concurrently with a
standalone DEFLATE implementation. This requires knowing, for every chunk,
where it lives in the file and how many bytes it occupies---that is, the
dataset's chunk index.

Obtaining that index is where a substantial and easily overlooked cost hides.
The natural approach queries the index once per chunk, but each such query
re-walks the dataset's chunk index from the beginning, so the total cost grows
quadratically with the number of chunks rather than linearly. We measured this
directly on GOES-19 imagery: a variable with $2304$ chunks took $0.042$~s to
index, while one with $9216$ chunks took $0.676$~s---four times the chunks for
sixteen times the time. At $0.5$~km this consumed more time than the
decompression it was meant to enable. Retrieving the entire index in a single
traversal, available from HDF5 1.14 onward, removes the quadratic term
altogether; the per-chunk fetch that follows remains linear at roughly
$20\,\mu$s per chunk and is not a concern. We note this because the quadratic
behavior is invisible in profiles taken at coarser resolutions, where the chunk
count is small enough that the index cost disappears into the noise: it emerges
precisely at the full-disk, full-resolution scale that operational use demands.

A second constraint applies to the fetch itself. The HDF5 library serializes
access through a global lock, so reading chunks through it cannot proceed in
parallel regardless of how many threads are available. Because the index
traversal already yields each chunk's byte offset within the file, the chunk
payloads can instead be read with \texttt{pread}, which takes an explicit
offset and does not share the file descriptor's position, and is therefore safe
to call concurrently from many threads. Two details matter for correctness:
offsets are relative to the file's base address, which is nonzero when a user
block is present and must be added; and chunks that the index reports with zero
stored size are unallocated regions that carry the fill value rather than data,
and must be materialized separately instead of being read. Where any
precondition fails---an older HDF5, an unavailable offset---the implementation
falls back to the library path, which is slower but correct.

The resulting design has a property worth stating plainly for anyone adopting
it: it does not depend on hpsatviews. It applies to any consumer of chunked,
DEFLATE-compressed NetCDF-4 that reads whole variables and can afford to
inflate them concurrently, which describes most operational satellite
processing.

\subsubsection{Device-resident processing}
\label{sec:device}

The per-pixel stages of a true-color composite---viewing geometry, atmospheric
correction, band synthesis, contrast transformation, compositing, and
reprojection---are individually well suited to a GPU. Offloading them one at a
time is nonetheless counterproductive: at full-disk resolution each stage
operates on arrays of order $10^{8}$ elements, so a kernel invoked in isolation
spends more time moving its operands across the host--device boundary than
computing on them. We confirmed this early in development, where a
straightforward port of a single per-pixel stage ran slower than the
multithreaded CPU implementation it replaced.

The backend is therefore organized around residency rather than around kernels.
Each input channel is transferred to the device once, and the entire chain is
executed there: the latitude/longitude grid is generated on the device from the
projection parameters rather than computed on the host and uploaded; solar and
satellite geometry follow from it; the lookup-table Rayleigh correction, the
synthetic green band, the ratio sharpening, the piecewise contrast stretch, and
the multiband composition each consume the previous stage's output in place.
The composed image itself remains resident, so when a geographic reprojection
is requested it reads the device copy directly. A full render therefore moves
the input channels in and one image out, with no intermediate round trips. The
transfer cost is paid once per scene instead of once per stage, which is what
makes the offload worthwhile at all.

Two consequences of this design bear on the results that follow. First,
generating the navigation grid on the device and computing viewing geometry
from it are double-precision--intensive: geolocation on the fixed grid, and the
inverse mapping performed by the reprojection, lose accuracy unacceptably in
single precision. The throughput of a given GPU in double precision therefore
governs whether the offload pays, and that ratio varies by more than an order
of magnitude across devices in the same product generation. This is the
mechanism behind the device-dependent verdict examined in
Section~\ref{sec:discussion}, and it is why a speed-up measured on one
accelerator should not be quoted for another.

Second, the CPU implementation remains the reference. Any option for which no
kernel exists falls back to it transparently, and the two paths are compared
against each other by the regression suite on every change. Retaining a
complete, independent implementation is not redundancy: it is what makes the
equivalence check of Section~\ref{sec:equivalence} possible, and it is what
allows a deployment to decline the GPU without losing functionality when
measurement says the GPU does not pay.

\subsubsection{Output encoding}
\label{sec:output}

Once decoding and computation are both parallel, the encoder becomes visible in
the profile, and two choices there matter more than they appear to.

The first concerns overviews. A Cloud-Optimized GeoTIFF carries an internal
pyramid of reduced-resolution copies, which is what allows a client to read a
coarse view over HTTP without retrieving the whole file. Building that pyramid
accounts for roughly $90\%$ of the time spent writing a full-disk GeoTIFF. In a
distribution pipeline, however, the georeferenced full disk is frequently an
intermediate: a downstream stage crops it to regional subsets, and the pyramid
is discarded without ever being read. Emitting a tiled but non-pyramidal
GeoTIFF by default, and constructing the full Cloud-Optimized form only when
the file is the delivered product, removes that cost from the common case
without removing the capability.

The second concerns how the pixels reach the encoder. Image data is held
interleaved by pixel, whereas a GDAL dataset is conventionally populated band
by band; the obvious implementation allocates a second image-sized buffer and
de-interleaves into it, a pass whose access pattern defeats the cache. Wrapping
the existing interleaved buffer in an in-memory dataset instead eliminates both
the allocation and the pass.

Finally, the choice of container is itself a performance decision at this
scale. The PNG encoder compresses on a single thread, and no configuration of
its filter and compression parameters changes that; for a full disk this makes
PNG about an order of magnitude more expensive to write than the equivalent
GeoTIFF. PNG remains appropriate for regional subsets, which is where it is
used, but a full-disk product written as PNG spends more time in the encoder
than in everything else combined.

%%%%%%%%%%%%%%%%%%%%%%%%%%%%%%%%%%%%%%%%%%%%%%%%%%%%%%%%%%%%%%%%%%%%%
\subsection{Testing: evaluation protocol}
\label{sec:methodology}

Performance claims about processing software are unusually easy to get wrong in
ways that survive review, because the resulting numbers are internally
consistent and reproducible---they simply measure something other than what
they are said to measure. In the course of this work we encountered three such
errors, each of which produced a plausible figure that we initially believed.
We describe them because the corrective practices generalize to any published
comparison of this kind, and because two of the three would have inflated our
own reported advantage.

\textit{Comparing builds from different revisions.} The natural way to quantify
a GPU backend is to time the accelerated build against the conventional one.
If the two timings are taken at different points in a development cycle,
however, the comparison silently absorbs every unrelated improvement made in
between. Over this cycle the CPU build alone improved from $3.76$~s to
$2.53$~s on the full-disk product, entirely from the decoding and encoding work
of Section~\ref{sec:implementation}, which benefits both builds equally.
Quoting the earlier CPU figure against the current GPU figure would have
reported $3.21\times$ where the honest value is $2.16\times$---a $49\%$
overstatement, crediting the accelerator with a second of work performed by the
input path. Both builds must be compiled from the same revision and timed in
the same session.

\textit{A silent fallback.} Accelerated implementations typically cover a
subset of the available options and revert to the reference implementation
otherwise. If that reversion is silent, a measurement can be numerically
correct and still be false about its cause. In our case the ratio-sharpening
option disqualified the composite from the accelerated path; because that same
option is required to match the product emitted by the tool we compare against
(Section~\ref{sec:results}), every cross-tool measurement we made was executing
the CPU path while reporting it as GPU. The wall times were real, but the
attribution was not. The correction was twofold: the tool emits an explicit
warning whenever a requested configuration falls back, and the regression suite
now asserts on the \emph{absence} of that warning for the configuration used in
the benchmark. This second point is the transferable one---an equivalence test
alone cannot detect this class of error, because both paths continue to produce
identical output. The fallback must be made loud, and the loudness itself must
be tested.

\textit{Comparing against another tool's defaults.} A comparison is only fair
if the reference tool is given a reasonable opportunity to perform. The tool we
compare against distributes work through a task scheduler whose degree of
parallelism is a user-set parameter, defaulting to four workers irrespective of
the host. On our 64-thread server it requires $108.3$~s at that default and
$30.6$~s at its best setting. Quoting the default would have reported an
$11.7\times$ advantage where the defensible value is $3.31\times$: an inflation
of $3.5\times$, obtained without misreporting a single measurement. We
therefore sweep that parameter and quote the reference tool's best time.

The protocol adopted from these lessons is as follows. Both builds are compiled
from the same source revision and timed in the same session on the same host.
The input files are read once before timing so that all configurations run
against a warm page cache and the comparison does not degenerate into a
measurement of storage. Each configuration is repeated and the median reported.
The comparison tool's parallelism parameter is swept and its best result used.
Every configuration is verified to produce the intended product, both by
checking that no fallback warning was emitted and by comparing the output
against the reference implementation. The scripts implementing this protocol
are distributed with the software, so the tables that follow can be regenerated
on other hardware.

\subsubsection{Experimental setup}
\label{sec:setup}

All timings were obtained on the laboratory's production ingest server: two
Intel Xeon Gold 6326 processors at $2.90$~GHz ($32$ cores, $64$ threads) and one
NVIDIA A30 accelerator. The scene is a GOES-19 ABI full disk. The measured
product is a true-color composite with lookup-table Rayleigh correction,
ratio sharpening, and the piecewise contrast stretch, rendered at the
$0.5$~km native resolution of the sharpened composite
($21{,}696 \times 21{,}696$ pixels) and written as a tiled GeoTIFF. This
configuration was chosen because it is the product the reference tool emits, so
the two engines can be compared on identical output rather than on their
respective defaults.

Reported times are end-to-end wall times of the complete process, including
startup, context creation, reading, and writing---not the duration of an
instrumented inner region. This is deliberate: in an operational pipeline the
quantity that matters is the interval from the arrival of a file to the
availability of an image, and startup costs that a profiler would exclude are
part of that interval.

%%%%%%%%%%%%%%%%%%%%%%%%%%%%%%%%%%%%%%%%%%%%%%%%%%%%%%%%%%%%%%%%%%%%%
\section{Results}
\label{sec:results}

\subsection{Per-scene wall time}
\label{sec:walltime}

Table~\ref{tab:sweep} reports the reference tool's wall time as a function of
its worker count. Throughput improves steeply up to sixteen workers, gains
little between sixteen and thirty-two, and degrades beyond that: on this
64-thread host the scheduler's own coordination overhead outweighs the
additional concurrency. The minimum, $30.63$~s at thirty-two workers, is the
figure we carry forward. We stress that this is more than three times faster
than the same tool at its shipped default, and that a comparison omitting this
sweep would not have been wrong in its arithmetic---only in its conclusion.

\begin{table}[t]
\caption{Wall time of the reference tool (geo2grid 1.3) on the full-disk
product as a function of its worker count, on the host of
Section~\ref{sec:setup}. Median of two runs.}\label{tab:sweep}
\centering
\begin{tabular}{lrrrrr}
\toprule
Workers & 4 & 8 & 16 & 32 & 64 \\
\midrule
Wall time (s) & 108.3 & 61.2 & 39.0 & \textbf{30.6} & 33.1 \\
\bottomrule
\end{tabular}
\end{table}

Table~\ref{tab:walltime} sets that figure against the two execution paths of
this work. Two comparisons are present and should not be conflated.

The first is like for like: the CPU path against the reference tool, on the
same processors, over the same scene, producing the same product. The
$3.31\times$ reduction is attributable to the implementation substrate alone---a
compiled binary with parallel decoding and encoding, against a Python
scheduler over the same algorithms. This is the figure we regard as the
paper's central quantitative result, because it isolates the variable of
interest.

The second compares the GPU path against a tool that has none. The
$10.14\times$ is what an operator on this hardware actually obtains, and it is
the honest answer to ``how much sooner does the image appear,'' but it spans
two classes of hardware and should be read as such. Within this work the
accelerator contributes $3.06\times$ over the engine's own CPU path
($9.25 \rightarrow 3.02$~s). We note that this exceeds the $2.16\times$ we
measure on the same host for the $1$~km product: at four times the pixel count
the balance shifts from input and output toward arithmetic, which is the
regime in which offloading pays.

\begin{table}[t]
\caption{End-to-end wall time for the full-disk product of
Section~\ref{sec:setup}, including process startup, reading and writing.
Median of two runs with a warm page cache. Speed-up is relative to the
reference tool at its best worker count.}\label{tab:walltime}
\centering
\begin{tabular}{lrr}
\toprule
Engine & Wall time (s) & Speed-up \\
\midrule
geo2grid 1.3 (best of Table~\ref{tab:sweep}) & 30.63 & --- \\
This work, CPU path (64 threads) & 9.25 & $3.31\times$ \\
This work, GPU path (NVIDIA A30) & 3.02 & $10.14\times$ \\
\bottomrule
\end{tabular}
\end{table}

\subsection{Product equivalence}
\label{sec:equivalence}

A faster engine that produced a different image would not be a substitute, so
the timings above are meaningless without a statement of what the two engines
agree on. We compared the two full-disk outputs over the illuminated disk,
excluding the off-disk corners, which carry fill rather than data.

Agreement is close but not exact. The mean absolute difference is $3.90$ of
$255$ digital counts, or $1.5\%$ of full scale; $95.9\%$ of samples fall within
$10$ counts and $98.2\%$ within $25$. The distribution is heavy-tailed---the
root-mean-square difference, $20.1$ counts, greatly exceeds the mean absolute
difference---and the residual is concentrated at the limb and along the
day/night terminator, where the two implementations differ in how the
solar-zenith normalization is applied at extreme angles. Mean scene brightness
differs by $0.8$ counts ($118.4$ against $117.6$).

The per-channel biases are not symmetric. Green and blue agree to within $0.7$
counts, while the red channel of this work is consistently darker by $2.4$
counts. We have reproduced this offset on two hosts, two scenes, and two
versions of the reference tool, so it is systematic rather than incidental, and
we report it as an unresolved discrepancy rather than absorb it into the
aggregate statistics. Its magnitude is well below the threshold of visual
significance for the products concerned, but it indicates a genuine difference
in the treatment of the red band---most plausibly in the atmospheric correction
applied to it, which is the only stage acting on that channel alone---and it
should be resolved before the two engines are treated as interchangeable at the
radiometric level.

Between the CPU and GPU paths of this work the agreement is of a different
order: fewer than $10^{-5}$ of samples differ at all, and none by more than one
digital count, consistent with floating-point rounding rather than with any
difference in formulation.

%%%%%%%%%%%%%%%%%%%%%%%%%%%%%%%%%%%%%%%%%%%%%%%%%%%%%%%%%%%%%%%%%%%%%
\section{Discussion}
\label{sec:discussion}

\subsection{Where the remaining time is spent}
\label{sec:timebudget}

Summing the instrumented stages of the accelerated path on the A30 accounts for
$0.164$~s of arithmetic in a render whose total wall time is $1.17$~s. Roughly
one seventh of the interval is computation; the remainder is process startup,
reading and decoding the input, and encoding the output. Once the per-pixel
work has been moved to an accelerator, in other words, the problem has ceased
to be compute-bound, and further kernel optimization has little left to
recover. Groups considering this kind of work should expect the same inversion
and plan accordingly: after the first round of offloading, effort is better
spent on the input and output paths than on additional kernels. This is the
reasoning behind the decoding and encoding work of
Section~\ref{sec:implementation}, which was undertaken after the GPU backend
and which benefits both execution paths.

\subsection{Portability of the GPU verdict}
\label{sec:portability}

A speed-up measured on one accelerator is often reported as though it were a
property of the software. It is not. We ran the identical binary and the same
scene on a second host equipped with an NVIDIA T4---a card of the preceding
generation, widely deployed in institutional servers---and obtained a
qualitatively different answer, for two independent reasons.

The first is arithmetic precision. Geolocation on the satellite fixed grid, and
the viewing geometry derived from it, must be evaluated in double precision;
the remaining stages operate in single precision. The two device generations
differ enormously in their relative double-precision capability, and
Table~\ref{tab:stages} shows the consequence stage by stage. The
double-precision stages are $15.4\times$ slower on the T4, while the
single-precision stages are only $1.9\times$ slower. The single-precision
stages act as an internal control: they establish that the difference is not a
general disparity between the cards, but is localized precisely where double
precision is required, and its magnitude corresponds to the ratio of the two
devices' double-precision throughputs. The practical consequence is that these
two stages consume $78\%$ of the arithmetic on the T4 against $30\%$ on the
A30, so the same code has a different bottleneck on each.

\begin{table}[t]
\caption{Per-stage device time for the $1$~km full-disk product, identical
binary and scene on both accelerators. Stages are grouped by the precision in
which they are evaluated.}\label{tab:stages}
\centering
\begin{tabular}{lrrr}
\toprule
Stage & T4 (s) & A30 (s) & Ratio \\
\midrule
\multicolumn{4}{l}{\textit{Double precision}} \\
\quad Latitude/longitude grid       & 0.153 & 0.015 & $10.2\times$ \\
\quad Solar and satellite geometry  & 0.616 & 0.035 & $17.6\times$ \\
\quad \textit{subtotal}             & 0.769 & 0.050 & $15.4\times$ \\
\midrule
\multicolumn{4}{l}{\textit{Single precision}} \\
\quad Solar zenith correction ($\times 3$) & 0.021 & 0.006 & $3.5\times$ \\
\quad Rayleigh lookup (2 channels)  & 0.021 & 0.011 & $1.9\times$ \\
\quad Synthetic green               & 0.009 & 0.003 & $3.0\times$ \\
\quad Piecewise stretch ($\times 3$) & 0.015 & 0.009 & $1.7\times$ \\
\quad Composition and readback      & 0.147 & 0.085 & $1.7\times$ \\
\quad \textit{subtotal}             & 0.213 & 0.114 & $1.9\times$ \\
\midrule
Total device arithmetic             & 0.982 & 0.164 & $6.0\times$ \\
\bottomrule
\end{tabular}
\end{table}

The second reason is memory capacity, and it is a direct cost of the residency
strategy that makes the offload worthwhile in the first place. Holding the
whole chain on the device means holding, simultaneously, three input channels,
two coordinate grids and three geometry grids. At the $0.5$~km resolution of
the sharpened product each such array occupies $1.9$~GB, for a working set of
approximately $15$~GB. On the A30 this fits; on the $16$~GB T4 the allocation
fails and the render reverts to the CPU path in its entirety. The product used
for the cross-tool comparison of Section~\ref{sec:results} therefore cannot be
GPU-accelerated at all on that card, while the same product at $1$~km, whose
working set is a quarter of the size, runs entirely on the device.

Neither limitation is a defect of the accelerator or of the design; both are
consequences of a specific pairing of algorithm and hardware. What they
establish is that a deployment cannot infer its own performance from a
published figure. The measurement must be repeated per host, which is why the
benchmark scripts are distributed with the software and why we report a
protocol rather than only a number.

\subsection{A systematic radiometric offset}
\label{sec:redbias}

The residual red-channel offset reported in Section~\ref{sec:equivalence}
remains unexplained. Its reproducibility across hosts, scenes and versions of
the reference tool rules out an incidental cause, and its confinement to a
single channel points to a stage acting on that channel alone---most plausibly
the atmospheric correction, since the red band is the one for which the
correction is evaluated against the reflectance of the same band. We report the
offset rather than absorb it because a difference of this kind, though below
the threshold of visual significance, would matter to any quantitative use of
the output, and because its resolution is a prerequisite for treating the two
implementations as radiometrically interchangeable.

\subsection{Limitations}
\label{sec:limitations}

The evaluation covers a single instrument, and although the decoding strategy
of Section~\ref{sec:decoding} applies to chunked NetCDF-4 generally, its
benefit depends on chunk geometry and compression settings that vary between
data producers. Timings come from two hosts, which is enough to demonstrate
that the GPU verdict is host-dependent but not to characterize the space of
devices. The engine targets a single multi-core node by design and has not been
evaluated in a distributed configuration. Finally, the comparison addresses one
product; other composites share the same decoding and encoding paths but
differ in their arithmetic, and their balance between the two would have to be
measured separately.

%%%%%%%%%%%%%%%%%%%%%%%%%%%%%%%%%%%%%%%%%%%%%%%%%%%%%%%%%%%%%%%%%%%%%
\section{Conclusions}
\label{sec:conclusions}

For a standard GOES-R composite at operational resolution, most of the
per-scene time is recoverable without touching the algorithms that define the
product. Decoding the input by retrieving raw chunks and inflating them
concurrently, encoding without an overview pyramid the downstream stages will
discard, and moving the per-pixel arithmetic to an accelerator as one resident
chain rather than as isolated kernels, together reduce the wall time for a
full-disk true-color composite from $30.6$~s to $9.25$~s on the same processors
as the reference implementation, and to $3.02$~s when a suitable GPU is
available. The outputs agree to a mean absolute difference of $1.5\%$ of full
scale.

Two results qualify that headline in ways we consider more useful than the
figure itself. First, the cost of obtaining the HDF5 chunk index grows
quadratically with the chunk count when queried per chunk; at $0.5$~km this
exceeded the decompression it was meant to enable, and it is undetectable in
profiles taken at coarser resolution. Second, whether GPU offloading pays is a
property of the pairing between this workload and a particular device, not of
the approach: the geolocation stages require double precision and are
$15.4\times$ slower on a T4 than on an A30, while the single-precision stages
differ by only $1.9\times$, and at full resolution the residency strategy
exceeds the T4's memory entirely, so the accelerated path is simply
unavailable there.

Both findings point the same way. A deployment cannot adopt a published
speed-up; it must measure its own host, and the measurement is easy to get
wrong in ways that survive scrutiny. We have therefore reported the protocol
alongside the numbers and distributed the scripts that implement it, so that
the tables presented here can be regenerated rather than believed. With the
arithmetic on an accelerator, the remaining wall time is dominated by input and
output rather than by computation, which is where we expect further work in
this class of pipeline to be worthwhile.


%%%%%%%%%%%%%%%%%%%%%%%%%%%%%%%%%%%%%%%%%%%%%%%%%%%%%%%%%%%%%%%%%%%%%
\section*{List of abbreviations}
%% TODO completar. Springer lo recomienda y es barato. Al menos: ABI, COG, DN,
%% FP32, FP64, GOES-R, HDF5, LUT, SZA, VZA.

%%%%%%%%%%%%%%%%%%%%%%%%%%%%%%%%%%%%%%%%%%%%%%%%%%%%%%%%%%%%%%%%%%%%%
\section*{Statements and Declarations}
%% OBLIGATORIO: las solicitudes que no incluyan las declaraciones pertinentes
%% se devuelven como incompletas. Van antes de la lista de referencias.
%%
%% OJO: Author contributions y Competing interests hay que capturarlas ADEMAS
%% en la interfaz de envio, y solo lo capturado ahi sale en la version
%% publicada. Ponerlas nada mas aqui no basta.

\subsection*{Funding}
This work was developed at the Laboratorio Nacional de Observaci\'on de la
Tierra (LANOT), Universidad Nacional Aut\'onoma de M\'exico, with support from
Laboratorio Nacional SECIHTI 2025--2027, grant LN-2025-C-102.

\subsection*{Competing interests}
%% TODO redactar.

\subsection*{Data availability}
%% TODO redactar: GOES-R ABI L1b del bucket publico de NOAA en S3, accesible
%% sin credenciales.

\subsection*{Code availability}
%% TODO redactar:
%%  - hpsatviews v1.1.0 archivado, DOI 10.5281/zenodo.21893553
%%    (DOI de concepto 10.5281/zenodo.20817973)
%%  - reproduction/bench_geo2grid.sh y reproduction/compare_g2g_product.sh
%%    reproducen las tablas de la seccion de resultados.

\subsection*{Author contributions}
%% TODO en formato CRediT. BLOQUEADO por la decision de coautores.

\subsection*{Use of large language models}
%% TODO decidir. La revista exige documentar el uso de un LLM y exime solo la
%% "edicion de copia asistida" ---legibilidad, estilo, gramatica, tono---,
%% excluyendo expresamente el trabajo editorial generativo y la creacion
%% autonoma de contenido. Ver la seccion 5 de paper/plan_esin.md.

%%%%%%%%%%%%%%%%%%%%%%%%%%%%%%%%%%%%%%%%%%%%%%%%%%%%%%%%%%%%%%%%%%%%%
%% plainnat mientras se trabaja con `article`. Con sn-jnl se cambia al estilo
%% que trae la plantilla.
\bibliographystyle{plainnat}
\bibliography{paper}

\end{document}
